%%%%%%%%%%%%%%%%%%%%%%%%%%%%%%%%%%%%%%%%%%%%%%%%%%%%%%%%%%%%%%%%%%%%%%%%%%%%%%%%
%2345678901234567890123456789012345678901234567890123456789012345678901234567890
%        1         2         3         4         5         6         7         8

\documentclass[letterpaper, 10 pt, conference]{ieeeconf}  % Comment this line out if you need a4paper

%\documentclass[a4paper, 10pt, conference]{ieeeconf}      % Use this line for a4 paper

\IEEEoverridecommandlockouts                              % This command is only needed if
                                                          % you want to use the \thanks command

\overrideIEEEmargins                                      % Needed to meet printer requirements.

%In case you encounter the following error:
%Error 1010 The PDF file may be corrupt (unable to open PDF file) OR
%Error 1000 An error occurred while parsing a contents stream. Unable to analyze the PDF file.
%This is a known problem with pdfLaTeX conversion filter. The file cannot be opened with acrobat reader
%Please use one of the alternatives below to circumvent this error by uncommenting one or the other
%\pdfobjcompresslevel=0
%\pdfminorversion=4

% See the \addtolength command later in the file to balance the column lengths
% on the last page of the document

% The following packages can be found on http://www.ctan.org
%\usepackage{graphics} % for pdf, bitmapped graphics files
%\usepackage{epsfig} % for postscript graphics files
%\usepackage{mathptmx} % assumes new font selection scheme installed
%\usepackage{times} % assumes new font selection scheme installed
%\usepackage{amsmath} % assumes amsmath package installed
%\usepackage{amssymb}  % assumes amsmath package installed

\usepackage{graphicx}
\usepackage{xcolor}
\newcommand{\TBD}[1]{\textbf{\color{red}[TBD-#1]}}

\title{\LARGE \bf XXXXX

}


\author{
}


\begin{document}



\maketitle
\thispagestyle{empty}
\pagestyle{empty}


%%%%%%%%%%%%%%%%%%%%%%%%%%%%%%%%%%%%%%%%%%%%%%%%%%%%%%%%%%%%%%%%%%%%%%%%%%%%%%%%
% ================================================================
% ABSTRACT(回填首页)
% ================================================================
\begin{abstract}
We present a multimodal command interface in which speech and gaze
jointly control a quadruped mobile manipulator. A head-worn eye
tracker and the robot share a persistent, metric instance map built
on 3D Gaussian Splatting; the wearer's gaze is cast into this shared
world model as a query over labeled object instances, so spoken
deictic references are grounded directly in 3D space rather than in
any camera image. Gaze thus resolves the ambiguity that language
alone cannot: the referent may be one of many identical objects,
unmarked, and outside the robot's line of sight. Resolution combines
world-frame fixation clustering, cone-posterior voting over labeled
Gaussians with online accuracy stamps, an eye--voice binding window,
and last-meter re-projection into the grasp camera. Across 157
pre-registered trials, top-1 resolution accuracy is 98.6\% at
angular separations above 2.5$^\circ$ --- identical, unmarked
instances included --- and degrades where the calibrated
1$^\circ$ gaze cone predicts, approaching chance below
1$^\circ$; end-to-end fetch succeeds in
\TBD{E2-headline} of \TBD{E2-K} trials, with success statistically
indistinguishable whether or not the target is in the robot's view.
\end{abstract}


%%%%%%%%%%%%%%%%%%%%%%%%%%%%%%%%%%%%%%%%%%%%%%%%%%%%%%%%%%%%%%%%%%%%%%%%%%%%%%%%
\section{INTRODUCTION}

Fetch-and-carry is the canonical service task of a mobile manipulator,
and speech is its natural interface. But ask a robot to fetch ``the
tennis ball'' when three identical ones lie on the table, and language
hits a structural wall: description --- class, color, size --- cannot
separate instances that satisfy every description equally. The request
is a one-in-$N$ lottery, for any channel that refers by properties.
Picking one requires an \emph{index} --- \emph{which one}, not
\emph{what kind} --- and humans already have one: when a speaker says
``\emph{that} one,'' their eyes are on the referent around the word
itself (the eye--voice span~\cite{eyevoicespan}).

Robot systems accept this index only inside their own viewport.
Existing gaze-commanded manipulation resolves reference in the robot's
current camera --- searching the live view for the attended
object~\cite{famhri,gamma2026,rueppel2025} --- or instruments objects
with fiducial tags~\cite{ibotassistant}.
Either way, what can be referred to collapses to what the robot
perceives right now. Human deictic reference does not work this way:
``that one'' designates an individual in the shared world, whether or
not any robot is looking. Reference is world-centric; the systems are
robot-centric.

We restore the venue. The room is reconstructed once as a 3D Gaussian
Splatting (3DGS) map, segmented into instances, and named; the wearer's
eye tracker and the robot localize against the same surveyed fiducials,
so a fixation becomes a query over \emph{world instances} --- not over
anybody's image --- and the robot (a Unitree B2 quadruped with a Z1
arm) enters, grasps, and delivers on the answer in the same coordinates
(Fig.~1). 3DGS is the representation because the venue demands four
things at once --- continuous geometry for $1^\circ$-precision ray
queries, instance-level structure, semantic association of 2D evidence
to persistent 3D labels, and persistent online querying; Sec.~III
grounds this choice against point clouds, meshes, and implicit fields.
Two consequences follow, and we test both: the referent needs to be in
no camera's view --- in our experiments the robot starts outside the
room, behind a doorway, with no line of sight to the target (E3) ---
and unmarked identical instances become separable in principle ---
gaze supplies the index that description cannot (E1).

The claim has a quantitative signature. If resolution is a geometric
query over the world model, accuracy among identical instances is
governed by the \emph{angular separation} of the competitors as seen
from the user, against the calibrated gaze noise
($\sigma\!\approx\!1^\circ$: about 5\,cm at 3\,m, commensurate with
balls 0.25\,m apart). We therefore measure accuracy as a function of
angular separation and ask whether it falls where $\sigma$ predicts
(E1, Fig.~4). Holding the $1^\circ$ budget end-to-end is the system's
engineering substance; Sec.~IV introduces each mechanism by the
measured failure of the naive design it replaced.

Our contributions are:
\begin{itemize}
\item a \textbf{world-centric resolution architecture}: one
persistent metric 3DGS instance map, mapped and named once, that a
head-worn eye tracker queries and a quadruped executes in --- so the
referent may be unmarked, identical to its neighbors, and outside
every camera's view (Sec.~III; tested in E1, E3);
\item the \textbf{pipeline that holds the $1^\circ$ gaze budget
end-to-end}: world-frame fixation clustering, cone-posterior voting
over labeled Gaussians, an eye--voice binding window, and last-meter
re-projection into the grasp camera --- each mechanism introduced by
the measured failure of the naive design it replaced (Sec.~IV;
ablated bit-exactly in E4);
\item a \textbf{falsifiable geometric account of gaze
disambiguation}: accuracy among identical instances as a function of
their angular separation, with the calibrated cone width predicting
where it falls --- established on a pre-registered real-robot
protocol, all runs recorded, none excluded (E1--E2, Fig.~4).
\end{itemize}

\section{RELATED WORK}

\subsection{Gaze-Guided Manipulation: the Index in a Viewport}
Multimodal deixis for machines dates to
Put-That-There~\cite{bolt1980}, and gaze has long served as a readout
of manipulation intent~\cite{belardinelli2024,huang2015} and a
control channel for assistive arms~\cite{fischerjanzen2024}. Its
current robotic form pairs a wearable tracker with speech or a
foundation model and resolves the referent in a \emph{current image}:
MR-headset gaze over a shared table with clarification
dialogue~\cite{rosen2020}; headset gaze calibrated into the robot's
frame to cluster the live point
cloud~\cite{weber2022,weber2023}; glasses gaze matched into the
robot's camera by features~\cite{famhri} or grounded by a
VLM~\cite{gamma2026}; robot-explored tabletop
graphs~\cite{glancesay}; egocentric keyframes~\cite{edith}; live
feeds from distributed cameras~\cite{rueppel2025}; class matching
from a headset ray~\cite{raycastgrasp2025}; semantic gaze histories
over tagged objects~\cite{semanticscanpath}. Speech-plus-gaze can
even disambiguate duplicates --- within the robot's 2D
view~\cite{humanoids2023}; gaze-conditioned policies handle lookalike
targets --- within the robot's current
observation~\cite{gaze2act2025}; and when the referent is \emph{not}
in view, the remedy on offer is to move the camera until it
is~\cite{yuguchi2019}. Across all of these, the referable world is
the currently perceived world: the human's index is accepted only
inside a viewport. We accept it in a persistent world model, so the
referent need not be in the \emph{robot's} view at all.

\subsection{Map-Mediated Mobile Manipulation: the Human Outside the Map}
A second line equips the robot with persistent semantic memory and
fetches from it: CLIP voxel maps scanned by phone~\cite{okrobot2024},
online-updated point-feature memories~\cite{dynamem2024},
open-vocabulary scene graphs~\cite{dovsg}, instance-level Gaussian
scene graphs that relocalize and update across
sessions~\cite{dgsgmind}, task-relevant Gaussians gathered during
search~\cite{helios}, robot-built language-embedded
splats~\cite{legs}, multi-robot semantic splats that even ingest
eyewear streams --- with gaze unused~\cite{hammer}, and full-home
semantic maps queried by voice for quadruped
fetch~\cite{beyondcybathlon}. These systems establish map-mediated
fetch --- for \emph{descriptions}. The query is language, or a hand:
drag-and-drop in a scanned dollhouse~\cite{holospot}, skeletal
pointing read by the robot's own camera against a prior
map~\cite{miel}. The human is never a \emph{resident} of the map ---
absent from its frame entirely, or a transient skeleton in the
robot's view --- and descriptive queries inherit the one-in-$N$ wall:
OK-Robot reports lookalike instances as a standing failure mode, with
user confirmation as the fallback~\cite{okrobot2024}; clarification
dialogue, the generic remedy, can only ask for more properties ---
exactly what identical instances lack~\cite{ingress2020,hatori2018}.
We keep the persistent map and change who inhabits it: the wearer is
localized in the same metric frame as the robot, and the query is the
gaze index itself.

\subsection{World-Frame Gaze: Measurement, Not Reference}
Casting wearable gaze into prebuilt reconstructions is established
--- as an instrument. Attention analytics localize the tracker and
intersect the ray: SLAM-localized gaze in RGB-D point
clouds~\cite{wang2018}, spatio-temporal gaze fields in dense
maps~\cite{attention4d}, semantic 3D fixation
mapping~\cite{semanticfovea}; a fixed arm has even answered
world-frame gaze, by pointing a laser at the fixated
surface~\cite{kogkas2017}. Egocentric datasets now ship every
ingredient side by side --- aligned trajectories, 3D gaze,
word-timestamped speech, multi-session splats --- composed into
nothing~\cite{aea2024}. These lines validate our geometry --- wearer
pose plus gaze ray against a metric model --- but not the referential
contract: no instance identity, no calibrated uncertainty, no actor
on the answering end resolving \emph{which one}. Going world-frame
also pays a methodological dividend: fixation detection under
locomotion otherwise needs inertial sensing and learned
classifiers~\cite{kothari2020}; in world coordinates it reduces to
clustering.

\subsection{Quantifying the Index: Selection Psychometrics and the Eye--Voice Span}
That gaze selection degrades in angular units is established on
screens and in headsets: success collapses psychometrically as
targets shrink toward tracker noise~\cite{schuetz2020}, calibrated
bivariate error models drive probabilistic selection among close
targets~\cite{evain2016,wei2023}, speech corrects gaze among crowded
ones~\cite{miniotas2006}, and Bayesian accumulation over candidates
is standard~\cite{bayesgaze2021} --- for labeled 2D targets,
distinguishable by content, with no world behind them. Natural
pointing, for comparison, resolves several degrees at
best~\cite{breuer2010}; wearable trackers hold $\sim$1$^\circ$ as
calibrated precision, not as free-motion
accuracy~\cite{onkhar2024}. On the temporal side, speakers fixate
the referent before naming it~\cite{eyevoicespan,meyer1998} --- a
lead that survives unconstrained egocentric
recording~\cite{lookandtell2025} --- and interfaces have bound gaze
to speech by fixed windows for two
decades~\cite{kaur2003,prasov2010,khan2026}, most recently by
word-timestamped fixation selection in a live-view robot
system~\cite{famhri}. Both halves are prior art we build on, not
claims we make. The composition is the claim: selection among
\emph{unmarked, physically identical} instances, in world
coordinates, with accuracy measured against angular separation and
predicted by an independently calibrated noise model (E1).

\subsection{Nearest Neighbors}
Two systems attack our conjunction most directly. iBotAssistant
couples wearable gaze and speech to a mobile manipulator that
retrieves and delivers --- but gaze returns an AprilTag identity, so
every referable object is instrumented, and localization rides on a
conventional navigation map rather than a shared instance
map~\cite{ibotassistant}. IntenBot realizes the closest information
flow --- headset deixis, a scene-graph memory, out-of-room targets
--- but its multi-room fetching is executed by a virtual robot, the
physical deployment is an armless base that only navigates, and gaze
is routinely supplemented by pointing~\cite{intenbot}. EgoSpot shares
our embodiment and shared-frame geometry, as mode-based teleoperation
toward a surface point rather than resolution of a spoken deictic to
an instance~\cite{egospot}. Perceptual anchoring has, separately,
fused gaze-wearable and robot streams into one object memory ---
without deixis or retrieval~\cite{sorensen2024}. To our knowledge, no
prior system resolves a spoken deictic by wearable gaze alone, in a
persistent metric instance map, to an unmarked target outside the
robot's view at command time, and then grasps and delivers it with a
mobile manipulator in one continuous physical trial --- and none, on
any substrate, has measured whether identical instances separate at
the rate calibrated gaze noise predicts.



%================================================================
%================================================================
\section{SYSTEM}
%================================================================
% Fig.2 三栏架构图在此引用;III-A 引 Fig.3 平面图。

A reference venue shared by a human and a robot (Fig.~2) puts four
requirements on the map. \emph{Continuous geometry}: a fixation is a
$\sim$1$^\circ$ cone, and ray queries at that precision tolerate
neither holes nor phantom surfaces. \emph{Instance-level structure}:
referents are individuals, so the primitives must group into nameable
instances. \emph{Semantic association}: 2D evidence must lift into
persistent 3D labels, so that gaze can vote per instance.
\emph{Persistent querying}: metric, localizable by both agents, and
fast enough to answer online --- a depth query renders in
\TBD{timing-depth} and a full cone verdict in \TBD{timing-cone}.
% 旧值:3ms / 4-6ms @TITAN X(宪法§8);建议在现工作站复测一次
3DGS meets all four in one representation: Gaussians render depth
along any ray (accumulated opacity doubling as an honest ``no
surface'' signal) while remaining discrete, labelable primitives that
masks lift onto and votes pool over. Point clouds offer primitives
without continuous coverage; meshes are brittle to lift labels onto;
implicit fields offer continuity without primitives.

\subsection{Mapping and the Instance Registry (once, offline)}
The room (Fig.~3) is photographed with a phone and reconstructed by
COLMAP with calibrated intrinsics locked (plane-dominant indoor
scenes leave a focal/depth gauge ambiguity that self-calibration
resolves poorly). A ChArUco board defines the metric world: board
corners are triangulated and a similarity transform aligns the
reconstruction to the \emph{board frame} (metric, $z$-up, floor at
$z\!\approx\!0$; residual \TBD{map-board-rms}).
% 旧值:v2 0.48mm;v6 重建后复测
A 3DGS model is trained with every pose-normalization step disabled,
so the splat lives in the board frame. Wall fiducials are surveyed
into the same frame by per-corner triangulation
(\TBD{map-tag-rms}); online localization later solves PnP
% 旧值:桌面 tag id79 2.2mm
directly against these surveyed corners, so \emph{physical tag size
never enters the online chain}. Fiducials are infrastructure for
localizing \emph{cameras} --- the wearer's and the robot's --- never
for identifying referents: no object in the registry carries a
marker.

Instances are lifted from 2D: SAM masks are back-projected through
rendered depth and snapped to Gaussians, turning each mask into a
Gaussian-ID set; masks agree across views by ID-set IoU (same-frame
edges are banned --- SAM's part/whole ladder would weld distinct
objects through transitive closure); containment-based merging and a
voxel split finish the instances. The discipline throughout:
\emph{2D proposes, 3D adjudicates}. A human names instances once;
naming merges (the name is the primary key across the whole stack).
The registry of referable objects is closed-world: named instances
plus background classes. A static interactability prior $\pi$
(one weight per name, class-uniform, fixed before capture; furniture
down-weighted) is defined here and used in Sec.~IV-C.
The map used in our experiments registers \TBD{map-images} images and
yields \TBD{map-instances} instances, \TBD{map-named} of them named;
the end-to-end geometric chain (calibration $\times$ survey $\times$
localization $\times$ map) verifies at \TBD{chain-deg}.
% 旧值:~0.1°(verify blend 0.93px);v6 复测

\subsection{Wearer State}
The tracker's fisheye world camera is localized per frame: ArUco
corners detected on the raw image, undistorted corner-wise, and
passed to planar PnP against the surveyed corners. A pose is
accepted through three gates (room and height bounds; mean
reprojection $\le$0.006 normalized $\approx$5\,px; jumps $>$1\,m
within 0.25\,s rejected, with forced re-acceptance after five
straight rejections) and interpolated across gaps up to 1\,s. The
head pose doubles as the user's position: every fixation carries it,
delivery anchors on it at confirmation time, and a pose older than
10\,s raises a staleness warning rather than a silent guess.

\subsection{Command Channel}
The command channel turns speech into structure and nothing more: an
LLM parses each transcribed utterance into an action $\in$ \{fetch,
go-to, stop\} and three reference slots (object / place /
destination), each filled by a \emph{name} or flagged as a
\emph{deictic} --- deictics are never guessed. Transcription is
continuous (open microphone, energy gate, every utterance archived
as audio). A deterministic hygiene layer scrubs deictic tokens that
leak into name fields and demotes verb-inconsistent actions. Parses
are cached; an entry persists only after the user confirms the
command it produced, so the cache stays human-audited. The emergency
stop word bypasses the LLM entirely (normalized word-table match).

\subsection{Robot Execution}
% ================= [此小节由狗端同学撰写] =====================
% 需覆盖:B2+Z1 硬件;狗端定位(同一套测绘 tag);导航;到位按名检测;
% 抓取;送达;急停实现。最后一米的"选哪只"规则在 IV-E(意图侧冻结),
% 这里写它的机器人侧实现(抓取相机、hint 经自身位姿投影、检测框匹配;
% 桌前位姿来源写清——新设计无桌面 tag)。
The robot consumes board-frame task orders: \texttt{target} = the
referent's centroid plus a suggested approach bearing (standoff and
obstacle avoidance are the robot's own), \texttt{object\_hint} = the
resolved instance centroid (Sec.~IV-E), and an optional delivery
point. It acknowledges within 100\,ms and reports progress on a
status stream. \TBD{robot-subsection}

\subsection{Design Discipline}
Resolution logic is deterministic end to end; the LLM sits in no fast
loop (parse-time only, cache-first); every event --- gaze verdicts,
commands, bindings, dispatches, robot statuses --- is appended to a
session log. Any live session replays bit-exactly, which is the basis
of the offline ablations (E4).


%================================================================
\section{REFERENCE RESOLUTION}
%================================================================
\emph{Notation.} $T_{WC}(t)=(R(t),\,\mathbf{o}(t))$ is the head pose
in the board frame; $\mathbf{g}(t)\in\mathbb{R}^2$ the gaze sample on
the undistorted normalized plane; $\beta(t)$ the gaze bias; $\sigma$
the calibrated cone width; $D(\mathbf{o},\mathbf{r})$ the splat depth
oracle along ray $\mathbf{r}$.

\subsection{World-Frame Fixation Events}
An image-frame fixation detector cannot see the fixations that matter
here: a walking user stabilizes gaze on a world point through the
vestibulo-ocular reflex, which reads as continuous saccading in the
image. We therefore detect fixations \emph{after} mapping each sample
into the world: $\mathbf{r}_W = R\,[\,\mathbf{g}-\beta_{\mathrm{eff}};
1\,]/\|\cdot\|$, $\;\mathbf{p} = \mathbf{o} +
D(\mathbf{o},\mathbf{r}_W)\,\mathbf{r}_W$ (median $z$-depth over the
central patch; low accumulated opacity returns \emph{no surface}).
Samples cluster incrementally with three angular rules: a point joins
while $\|\mathbf{p}-\mathbf{c}\| \le \max(0.05\,\mathrm{m},\,
\tan(1.5^\circ)\,d)$ with $d$ the viewing distance --- the radius
scales with distance exactly as gaze noise does; a split triggers
when apparent angular velocity exceeds $80^\circ$/s (a metric
threshold tuned near shreds fixations far, since noise-induced
apparent speed grows with $d$); a cluster closes after 0.4\,s without
a mappable sample and survives if it lasted $\ge$0.25\,s. Because a
fixation only closes when gaze \emph{leaves} the target, the still-open
cluster is judged provisionally every 0.4\,s; finals settle the books.

\subsection{Online Accuracy Stamps}
A \TBD{drift-deg} within-wear drift once invalidated an entire
% 旧值:2.6°(rec002,宪法§4.5);可复用旧案或在新图复测
session: constant-bias calibration dies quietly as the headset creeps.
Surveyed fiducials give millimeter ground truth for free, so bias is
re-measured \emph{during use}: an uninterrupted stare at one tag
(same tag, sample gaps $\le$0.1\,s, dwell $\ge$0.4\,s, $\ge$25
samples, MAD outlier rejection; both raw and corrected directions are
tested so a stale correction cannot lock out its own replacement)
yields a stamp $(\hat\beta, \hat\sigma)$; $\hat\sigma$ is the pooled
per-axis RMS of the cleaned residuals. Tags adjacent to experiment
objects are excluded as stamp targets. A stamp's authority decays
with age,
\begin{equation}
\beta_{\mathrm{eff}}(t) = \hat\beta\, e^{-(t - t_s)/\tau},
\qquad \tau = 45\,\mathrm{s},
\end{equation}
so an outdated correction fades toward raw gaze instead of pushing it
off target. One statistical rule governs everything downstream: a
fixation is \emph{one} observation, not $N$,
\begin{equation}
\sigma_{\mathrm{eff}}^2 = \sigma_{\mathrm{bias}}^2 +
\sigma_{\mathrm{jitter}}^2/N,
\end{equation}
because calibration error is shared across all samples of a fixation.
Treating samples as independent would sharpen the cone by $\sqrt{N}$
and produce confidently wrong verdicts.

\subsection{Cone-Posterior Instance Voting}
The naive resolver --- collect Gaussians in a sphere around the gaze
point --- is viewpoint-independent, which contradicts what a fixation
is: floor Gaussians \emph{behind} a prone object vote even though the
user cannot see them. On our validation recording, sphere voting
produced confident misjudgments (a wrong object carrying a
\TBD{sphere-fail-share} vote share); its replacement lets only
% 旧值:84%(次案 44%);rec002,宪法§4.5 口径按此修正
\emph{visible} surface vote. The patch renderer already computes an
$S\times S$ depth-and-opacity block ($S{=}33$, spanning $\pm2.5\sigma$);
each pixel casts a ray with weight
$w_i=\exp(-\theta_i^2/2\sigma^2)$, back-projects at its rendered
depth, and snaps to the nearest labeled Gaussian within 5\,cm
(contour pixels, whose blended depths float between objects, honestly
fail the snap). Raw votes and abstention are
\begin{equation}
v(\ell) = \textstyle\sum_{i:\,L(\mathbf{x}_i)=\ell} w_i \alpha_i,
\qquad
p_{\mathrm{none}} = 1 - \textstyle\sum_\ell v(\ell) \big/
\textstyle\sum_i w_i ,
\end{equation}
with the denominator over \emph{all} patch pixels --- invisible,
out-of-range, and unlabeled surface all counted as ``don't know.''
Votes pool by name (naming merges), and only then multiplies the
interactability prior: $\tilde V(n) = \pi(n)\sum_{\ell\in n}v(\ell)$.
Because $\pi$ is class-uniform, the argmax \emph{among same-class
instances} is invariant to it --- the prior suppresses cross-class
vote leakage only, which E4 ablates. The verdict published per
fixation is top-$k$ candidates with shares, $p_{\mathrm{none}}$, and
the $\sigma$ actually used --- a calibrated posterior, not a label.

\subsection{Eye--Voice Binding and Slot Resolution}
Each transcribed command is anchored at its utterance end $t_w$
(back-dated for transcription latency). Each deictic slot is bound to
gaze evidence in an asymmetric window: 4\,s back, 0.5\,s forward
grace, with a still-running fixation ranking first --- speakers look
at the referent before and while naming it (eye--voice span).
The \emph{object} slot binds to instance identity (verdict
$\ge$0.5 vote share, cone mode only; a noun class in the utterance
filters candidates; only \emph{named} instances are bindable). The
\emph{place} and \emph{destination} slots bind to the \emph{landing
point} --- the world point gaze actually struck ($\ge$0.4\,s dwell;
unnamed instances count; floor and walls do not) --- deliberately not
the instance centroid: ``this spot'' means the point you looked at.
Name-filled slots resolve by fuzzy match over the registry with a
unique-hit rule (accept only a clear single winner). Resolution runs
destination $\to$ place $\to$ object; any unresolved slot aborts with
a spoken refusal --- the system never substitutes a guess. A
confirmation gate precedes dispatch; the task order carries the
referent centroid, an approach bearing computed from the user's side,
and the instance centroid as \texttt{object\_hint}.

\subsection{Last-Meter Disambiguation}
Class-level detection at the table re-loses the instance: among $N$
identical balls, a detector's boxes are descriptively
indistinguishable --- the 1-in-$N$ wall returns at one meter. The
resolved instance's map centroid closes it: the robot projects the
\texttt{object\_hint} into its grasp camera through its own estimated
board-frame pose and picks the detection box containing the projected
point; if no box contains it and none is within half a box width, the
robot \emph{refuses} rather than guesses. No fiducial is needed at
the table: the projection rides on the same localization the robot
navigates by, and the short grasp-range standoff keeps the angular
lever arm small, so the projection error budget (robot pose at the
table $\times$ map centroid, \TBD{lastmeter-budget}) stays small
against the object radius and densely packed lookalikes remain
separable.
% 新设计(2026-08 demo 实测):无桌面 tag,hint 直接经狗自身位姿投影。
% 狗端同学确认三点:①桌前位姿来源(墙 tag / 里程计融合?)
% ②负例拒抓分支(hint_mismatch)在新实现里是否仍在——明天负例必测
% ③lastmeter-budget 改从 E2 日志实测:hint 投影点到目标框的残差中位数

\subsection{The Failure Envelope}
Every stage of this pipeline scales \emph{angularly}: the gaze noise
$\sigma$, the clustering radius, and the voting cone all subtend
fixed angles, so distance enters resolution only through the
\emph{angular separation} of competing instances,
$\Delta\theta \approx s/d$ for spacing $s$ at distance $d$. This
collapses ``how far'' and ``how dense'' onto one axis and yields a
prediction: accuracy among identical instances degrades to chance as
$\Delta\theta$ approaches $O(\sigma)$, i.e.\ beyond
$d^\ast \approx s/(k\sigma)$ for some $k$ of order unity. E1 measures
the curve and locates the knee (\TBD{E1-knee}); localization quality
and splat depth noise are logged per station so the knee is
attributable to the gaze channel rather than to a decaying pose.
Failures off this axis --- no fiducial in view, unmapped or moved
objects, a missed binding window, transcription errors --- do not
bend the curve; they are enumerated and counted in the E2 failure
taxonomy.


\section{EXPERIMENTS}

\subsection{Setup and Protocol}
The arena is two connected spaces separated by a doorway
(Fig.~3): the user and target objects occupy one room; the robot
parks in the other, without line of sight to the targets. Targets
include \TBD{n-balls} identical, unmarked tennis balls placed
\TBD{layout} apart. Ground truth follows a card protocol: each trial's
target is drawn from a pre-generated pseudorandom card sequence
(committed before collection); the physical layout and its
instance-ID correspondence are photo-archived; the system's choice is
the instance ID in the binding log. Metrics and the failure taxonomy
below were frozen before data collection; \emph{every trial is
recorded and reported --- none excluded}. All parameters of
Secs.~III--IV were frozen at commit \TBD{freeze-commit}.

\subsection{E1: Disambiguation Accuracy vs.\ Angular Separation}
E1 asks the core question --- when language cannot separate identical
instances, does gaze? --- without the robot in the loop. The user
fixates the carded target and issues ``bring me \emph{that} ball'';
commands are injected as text at logged timestamps, so ASR error is
excluded by construction (it is measured in E2). Independent
variables: instance count $N \in \{2,3,5\}$ (fixed layout, subset
selection) and angular separation, swept by crossing spacing
$\{0.25, 0.5, 1.0\}$\,m with standing distance $\{1.5, 3, 5, 7\}$\,m
--- the far, dense cells push $\Delta\theta$ into the
$1$--$2^\circ$ regime where Sec.~IV-G predicts collapse. Each cell
runs 15--20 trials (\TBD{E1-total} total). Per-station covariates
(fiducials in view, localization residual, online $\sigma$) are
logged so the knee is attributable to the gaze channel.

\emph{Baselines} (scored on the same trials): (i) language only ---
analytically $1/N$, plus an LLM given the utterance text and the
instance list; (ii) nearest same-class instance to the user.
\emph{Metrics}: top-1 accuracy with Wilson 95\% CIs per cell,
first-vs-second vote margin, and McNemar tests against both
baselines.

Fig.~4 plots accuracy against angular separation with the calibrated
$\sigma$ drawn as a vertical line. Accuracy is \TBD{E1-near} above
\TBD{E1-sep-safe} of separation and falls to \TBD{E1-far} as
separation approaches $\sigma$, with the knee at
\TBD{E1-knee}~--- consistent with the geometric prediction of
Sec.~IV-G. % 膝盖没测出来就如实报平线并讨论(但按设计 7m/0.25m 应该塌)
A walking condition repeats one column of the grid while the user
moves ($\ge$\TBD{E1-walk-n} trials): world-frame clustering holds
accuracy at \TBD{E1-walk} vs.\ \TBD{E1-stand} standing, isolating the
contribution of locomotion-stabilized fixation detection.

\subsection{E2: End-to-End Fetch (Table~I)}
K $\ge$ 25 trials across $\ge$2 object classes and 2--3 layouts, with
hard resets between trials (robot to park, objects re-placed). A
trial succeeds only with \emph{zero human intervention}. Table~I
reports staged success --- reference resolution / navigation to
within 0.3\,m of the approach pose / grasp / delivery --- and a
latency decomposition (utterance end $\to$ dispatch $\to$ arrival
$\to$ grasp $\to$ delivery). Every failure lands in exactly one
pre-registered class: localization gap / wrong-instance selection /
missed binding window / transcription error / navigation failure /
grasp failure / communication timeout.

\begin{table}[t]
\caption{E2 end-to-end results (staged success and latency).}
\label{tab:e2}
\centering
\begin{tabular}{lcc}
\hline
Stage & Success & Median latency (s)\\
\hline
Reference resolution & \TBD{}/\TBD{} & \TBD{} \\
Navigation ($<$0.3\,m) & \TBD{}/\TBD{} & \TBD{} \\
Grasp (correct instance) & \TBD{}/\TBD{} & \TBD{} \\
Delivery & \TBD{}/\TBD{} & \TBD{} \\
\hline
End-to-end & \TBD{}/\TBD{} & \TBD{} total \\
\hline
\end{tabular}
\end{table}

\subsection{E3: In-View / Out-of-View Equivalence}
The same fetch task under two command-time conditions: (a) the robot
in the same room, target inside its camera view; (b) the robot behind
the doorway, no line of sight --- 10--12 trials each. Map mediation
provides no mechanism by which (b) should be harder \emph{for
resolution}: we observe \TBD{E3-in} vs.\ \TBD{E3-out} success with
overlapping CIs and matching failure classes. The (a)/(b) contrast
also separates navigation difficulty from resolution ability. In-view
approaches (e.g.\ FAM-HRI-style matching) are structurally
inapplicable in (b); we state this by citation rather than
re-implementation.

\subsection{E4: Replay Ablations (Table~II)}
Every trial's event stream replays bit-exactly, so ablations cost no
robot time. Table~II sweeps: eye--voice window start/width (0--3\,s);
noun-class filtering on/off; online accuracy stamps on/off (stale
constant bias); vote-share threshold 0.3--0.7; interactability prior
on/off; candidate scope named-only vs.\ all instances.

\begin{table}[t]
\caption{E4 ablations: E1 accuracy under replayed variants.}
\label{tab:e4}
\centering
\begin{tabular}{lc}
\hline
Variant & Top-1 accuracy \\
\hline
Full system & \TBD{} \\
No online stamps (constant bias) & \TBD{} \\
No class filter & \TBD{} \\
Binding window \TBD{best}/\TBD{worst} & \TBD{}/\TBD{} \\
Prior off & \TBD{} \\
Candidate scope: all instances & \TBD{} \\
\hline
\end{tabular}
\end{table}
% 预期口径:prior off 不动同类消歧、只影响跨类漏票;scope=all 掉分
% 来自未命名碎片抢票——写结果时按这两句解读,别倒果为因。

\subsection{System Characterization}
From logs, no new collection: localization coverage \TBD{sys-loc},
post-stamp gaze $\sigma$ \TBD{sys-sigma}, depth query
\TBD{timing-depth}, cone verdict \TBD{timing-cone}, dispatch receipt
\TBD{sys-rep}, emergency-stop response \TBD{sys-estop} (measured
stop-request to motion-halt on hardware). Fig.~5 shows qualitative
strips: gaze cross, instance boxes, and verdict banners across
scenes.

%================================================================
\section{LIMITATIONS}
%================================================================
The map is static: moved objects invalidate both depth and centroids
until remapped (render-vs-photo blending flags changes but does not
repair them). Referability is closed-world twice over --- only named
instances can be bound, and the robot's detector vocabulary is a
finite class list. Calibration lives for one wearing; within-wear
drift is suppressed by the stamp protocol, which requires surveyed
fiducials in view from time to time, and head localization dies
without wall fiducials. Binding anchors at utterance end; a sentence
with two deictics relies on an ordering heuristic (word-level
timestamps are future work). Delivery targets the user's position at
confirmation time. Last-meter disambiguation inherits the robot's
localization error at the table; when the projection matches no
detection, the robot refuses rather than guesses. Transparent and unmapped surfaces
return no depth and cannot be referred to. Single user, single scene
(a cross-user study is future work).

%================================================================
\section{CONCLUSION}
%================================================================
Reference resolution moved venue: from the robot's camera to a
persistent metric 3DGS instance map that the human's gaze queries and
the robot acts in. The change dissolves two standing premises ---
in-view targets and instrumented objects --- and turns ``can gaze
tell identical instances apart'' into a geometric proposition with a
measurable failure point. A quadruped manipulator fetching an
unmarked, gaze-designated ball from outside the room is the
existence proof; the angular-separation curve is the science.



%%%%%%%%%%%%%%%%%%%%%%%%%%%%%%%%%%%%%%%%%%%%%%%%%%%%%%%%%%%%%%%%%%%%%%%%%%%%%%%%
\section*{APPENDIX}

Appendixes should appear before the acknowledgment.

\section*{ACKNOWLEDGMENT}

The preferred spelling of the word ÒacknowledgmentÓ in America is without an ÒeÓ after the ÒgÓ. Avoid the stilted expression, ÒOne of us (R. B. G.) thanks . . .Ó  Instead, try ÒR. B. G. thanksÓ. Put sponsor acknowledgments in the unnumbered footnote on the first page.



%%%%%%%%%%%%%%%%%%%%%%%%%%%%%%%%%%%%%%%%%%%%%%%%%%%%%%%%%%%%%%%%%%%%%%%%%%%%%%%%

References are important to the reader; therefore, each citation must be complete and correct. If at all possible, references should be commonly available publications.



\begin{thebibliography}{99}

% ---- \bibitem entries (IEEEtran style) for related_work_v2.tex ----
% Paste into your \begin{thebibliography}{99} ... \end{thebibliography}.
% Ordering: rearrange to match first-citation order in your final text.

% ===== New references introduced in RELATED WORK v2 =====

\bibitem{fischerjanzen2024}
A. Fischer-Janzen, T. M. Wendt, and K. Van Laerhoven, ``A scoping
review of gaze and eye tracking-based control methods for assistive
robotic arms,'' \emph{Front. Robot. AI}, vol. 11, p. 1326670, 2024.

\bibitem{weber2022}
D. Weber, E. Kasneci, and A. Zell, ``Exploiting augmented reality for
extrinsic robot calibration and eye-based human-robot collaboration,''
in \emph{Proc. ACM/IEEE Int. Conf. Human-Robot Interaction (HRI)},
2022, pp. 284--293.

\bibitem{weber2023}
D. Weber, W. Fuhl, E. Kasneci, and A. Zell, ``Multiperspective
teaching of unknown objects via shared-gaze-based multimodal
human-robot interaction,'' in \emph{Proc. ACM/IEEE Int. Conf.
Human-Robot Interaction (HRI)}, 2023, pp. 544--553.

\bibitem{yuguchi2019}
A. Yuguchi, T. Inoue, G. A. Garcia Ricardez, M. Ding, J. Takamatsu,
and T. Ogasawara, ``Real-time gazed object identification with a
variable point of view using a mobile service robot,'' in \emph{Proc.
IEEE Int. Conf. Robot Human Interact. Commun. (RO-MAN)}, 2019.

\bibitem{dgsgmind}
L. Ge, X. Zhu, J. Liu, and X. Li, ``DGSG-Mind: Dynamic 3D Gaussian
scene graphs for long-term scene understanding and grounding,''
arXiv:2605.29879, 2026.

\bibitem{ingress2020}
M. Shridhar, D. Mittal, and D. Hsu, ``INGRESS: Interactive visual
grounding of referring expressions,'' \emph{Int. J. Robot. Res.},
vol. 39, no. 2--3, pp. 217--232, 2020.

\bibitem{hatori2018}
J. Hatori \emph{et al.}, ``Interactively picking real-world objects
with unconstrained spoken language instructions,'' in \emph{Proc.
IEEE Int. Conf. Robot. Autom. (ICRA)}, 2018, pp. 3774--3781.

\bibitem{kogkas2017}
A. A. Kogkas, A. Darzi, and G. P. Mylonas, ``Gaze-contingent
perceptually enabled interactions in the operating theatre,''
\emph{Int. J. Comput. Assist. Radiol. Surg.}, vol. 12, no. 7,
pp. 1131--1140, 2017.

\bibitem{aea2024}
Z. Lv \emph{et al.}, ``Aria everyday activities dataset,''
arXiv:2402.13349, 2024.

\bibitem{schuetz2020}
I. Schuetz, T. S. Murdison, and M. Zannoli, ``A psychophysics-inspired
model of gaze selection performance,'' in \emph{Proc. ACM Symp. Eye
Tracking Res. Appl. (ETRA)}, 2020.

\bibitem{evain2016}
A. \'Evain, F. Argelaguet, G. Casiez, N. Roussel, and A. L\'ecuyer,
``Design and evaluation of fusion approach for combining brain and
gaze inputs for target selection,'' \emph{Front. Neurosci.}, vol. 10,
p. 454, 2016.

\bibitem{wei2023}
Y. Wei, R. Shi, D. Yu, Y. Wang, Y. Li, L. Yu, and H.-N. Liang,
``Predicting gaze-based target selection in augmented reality headsets
based on eye and head endpoint distributions,'' in \emph{Proc. CHI
Conf. Human Factors Comput. Syst.}, 2023.

\bibitem{miniotas2006}
D. Miniotas, O. \v{S}pakov, I. Tugoy, and I. S. MacKenzie,
``Speech-augmented eye gaze interaction with small closely spaced
targets,'' in \emph{Proc. Symp. Eye Tracking Res. Appl. (ETRA)},
2006, pp. 67--72.

\bibitem{bayesgaze2021}
Z. Li \emph{et al.}, ``BayesGaze: A Bayesian approach to eye-gaze
based target selection,'' in \emph{Proc. Graphics Interface}, 2021,
pp. 231--240.

\bibitem{breuer2010}
T. Breuer, P. G. Pl\"oger, and G. K. Kraetzschmar, ``Precise pointing
target recognition for human-robot interaction,'' in \emph{Proc.
SIMPAR Workshops}, 2010, pp. 229--240.

\bibitem{onkhar2024}
V. Onkhar, D. Dodou, and J. C. F. de Winter, ``Evaluating the Tobii
Pro Glasses 2 and 3 in static and dynamic conditions,'' \emph{Behav.
Res. Methods}, vol. 56, no. 5, pp. 4221--4238, 2024.

\bibitem{meyer1998}
A. S. Meyer, A. M. Sleiderink, and W. J. M. Levelt, ``Viewing and
naming objects: Eye movements during noun phrase production,''
\emph{Cognition}, vol. 66, no. 2, pp. B25--B33, 1998.

\bibitem{lookandtell2025}
A. Deichler and J. Beskow, ``Look and tell: A dataset for multimodal
grounding across egocentric and exocentric views,''
arXiv:2510.22672, 2025.

\bibitem{kaur2003}
M. Kaur \emph{et al.}, ``Where is `it'? Event synchronization in
gaze-speech input systems,'' in \emph{Proc. Int. Conf. Multimodal
Interfaces (ICMI)}, 2003, pp. 151--158.

\bibitem{prasov2010}
Z. Prasov and J. Y. Chai, ``Fusing eye gaze with speech recognition
hypotheses to resolve exophoric references in situated dialogue,'' in
\emph{Proc. Conf. Empirical Methods Natural Lang. Process. (EMNLP)},
2010, pp. 471--481.

\bibitem{khan2026}
A. A. Khan \emph{et al.}, ``Gaze and speech in multimodal
human-computer interaction: A scoping review,'' in \emph{Proc. CHI
Conf. Human Factors Comput. Syst.}, 2026.

\bibitem{sorensen2024}
S. L. S{\o}rensen, S. Huang, Y. Cao, and M. B. Kj{\ae}rgaard,
``Perceptual anchoring for gaze-tracking wearables and robot-mounted
sensors,'' in \emph{Proc. Int. Conf. Ubiquitous Robots (UR)}, 2024,
pp. 71--76.

% ===== Intro "% TODO bibitem" resolutions =====

\bibitem{eyevoicespan}
Z. M. Griffin and K. Bock, ``What the eyes say about speaking,''
\emph{Psychol. Sci.}, vol. 11, no. 4, pp. 274--279, 2000.

\bibitem{famhri}
Y. Lai, S. Yuan, P. Li, B. Zhang, B. Kiefer, T. Deng, and A. Zell,
``FAM-HRI: Foundation-model assisted multi-modal human-robot
interaction combining gaze and speech,'' \emph{IEEE Trans. Autom.
Sci. Eng.}, vol. 23, pp. 10521--10534, 2026.

\bibitem{gamma2026}
T. Y. H. Tay, X. Yan, J. Ouyang, D. Wu, W. Jiang, J. Kao, and Y. Cui,
``Intent at a glance: Gaze-guided robotic manipulation via foundation
models,'' in \emph{RSS Workshop on Robot Planning in the Era of
Foundation Models}, 2025.

\bibitem{rueppel2025}
J. V. R\"uppel, A. Rudenko, T. Schreiter, M. Magnusson, and
A. J. Lilienthal, ``Gaze-supported large language model framework for
bi-directional human-robot interaction,'' in \emph{Proc. IEEE Int.
Conf. Robot Human Interact. Commun. (RO-MAN)}, 2025, pp. 677--684.

% ===== Nearest neighbors (verified 2026-08-10; see tmp/related_work_scan/out/11-nearest-neighbors.md) =====

\bibitem{ibotassistant}
X. Liang and J. Cai, ``A multimodal agentic AI framework for intuitive
human--robot collaboration,'' \emph{Sensors}, vol. 26, no. 6,
p. 1958, 2026.

\bibitem{intenbot}
Y.-T. Liu \emph{et al.}, ``IntenBot: Flexible and imprecise multimodal
input for LLMs to understand user intentions for casual and
human-like HRI,'' arXiv:2605.04585, 2026.

\bibitem{egospot}
G. Zhang \emph{et al.}, ``EgoSpot: Egocentric multimodal control for
hands-free mobile manipulation,'' arXiv:2306.02393, 2023, rev. 2026.

% ===== 补缺 23 条(作者已核实 2026-08-10;验证表与来源 URL 见
%       tmp/related_work_scan/out/bib-authors-{A,B,C}.md)=====

\bibitem{bolt1980} R. A. Bolt, ```Put-that-there': Voice and gesture
at the graphics interface,'' in \emph{Proc. SIGGRAPH}, 1980.

\bibitem{belardinelli2024} A. Belardinelli, ``Gaze-based intention
estimation: principles, methodologies, and applications in HRI,''
\emph{ACM Trans. Human-Robot Interaction}, 2024.

\bibitem{huang2015} C.-M. Huang, S. Andrist, A. Saupp\'e, and
B. Mutlu, ``Using gaze patterns to predict task intent in
collaboration,'' \emph{Frontiers in Psychology}, vol.~6, 2015.

\bibitem{rosen2020} E. Rosen, D. Whitney, M. Fishman, D. Ullman, and
S. Tellex, ``Mixed reality as a bidirectional communication interface
for human-robot interaction,'' in \emph{Proc. IEEE/RSJ IROS}, 2020,
pp. 11431--11438.

\bibitem{glancesay} Y. Lai, S. Yuan, P. Li, B. Kiefer, and A. Zell,
``Glance-Say: Multimodal human-robot collaboration and intent
recognition via sticky glance,'' arXiv:2603.06121, 2026.

\bibitem{edith} D. Lee, J. Choi, D. K. Shin, S. Kang, and K. Lee,
``Hierarchical policies from verbal and egocentric human signals for
natural human-robot interaction,'' arXiv:2606.10276, 2026.

\bibitem{raycastgrasp2025} Z. Lin, Y. Sang, and Y. Ye,
``RaycastGrasp: Eye-gaze interaction with wearable devices for
robotic manipulation,'' arXiv:2510.22113, 2025.

\bibitem{semanticscanpath} E. Menendez, M. Gienger, S. Mart\'inez,
C. Balaguer, and A. Belardinelli, ``SemanticScanpath: Combining gaze
and speech for situated human-robot interaction using LLMs,''
arXiv:2503.16548, 2025.

\bibitem{humanoids2023} Y.-C. Chang, N. Gandi, K. Shin, Y.-J. Mun,
K. Driggs-Campbell, and J. Kim, ``Specifying target objects in robot
teleoperation using speech and natural eye gaze,'' in \emph{Proc.
IEEE-RAS Humanoids}, 2023, pp. 1--7.

\bibitem{gaze2act2025} K. Zuo \emph{et al.}, ``Gaze2Act:
Gaze-conditioned vision-language-action policies for interactive
robot manipulation,'' arXiv:2605.30282, 2026.

\bibitem{okrobot2024} P. Liu, Y. Orru, C. Paxton, N. M. M. Shafiullah,
and L. Pinto, ``OK-Robot: What really matters in integrating
open-knowledge models for robotics,'' in \emph{Proc. Robotics:
Science and Systems (RSS)}, 2024.

\bibitem{dynamem2024} P. Liu, Z. Guo, M. Warke, S. Chintala,
C. Paxton, N. M. M. Shafiullah, and L. Pinto, ``DynaMem: Online
dynamic spatio-semantic memory for open world mobile manipulation,''
in \emph{Proc. IEEE ICRA}, 2025, pp. 13346--13355.

\bibitem{dovsg} Z. Yan \emph{et al.}, ``Dynamic open-vocabulary 3D
scene graphs for long-term language-guided mobile manipulation,''
\emph{IEEE Robot. Autom. Lett.}, vol. 10, no. 5, pp. 4252--4259,
2025.

\bibitem{helios} K. Ashton \emph{et al.}, ``HELIOS: Hierarchical
exploration for language-grounded interaction in open scenes,''
arXiv:2509.22498, 2025.

\bibitem{legs} J. Yu \emph{et al.}, ``Language-embedded Gaussian
splats (LEGS): Incrementally building room-scale representations with
a mobile robot,'' in \emph{Proc. IEEE/RSJ IROS}, 2024,
pp. 13326--13332.

\bibitem{hammer} J. Yu, T. Chen, and M. Schwager, ``HAMMER:
Heterogeneous, multi-robot semantic Gaussian splatting,'' \emph{IEEE
Robot. Autom. Lett.}, vol. 10, no. 7, pp. 7270--7277, 2025.

\bibitem{beyondcybathlon} C. Scheidemann, A. Cramariuc, C. Chen,
J.-R. Chiu, and M. Hutter, ``Beyond Cybathlon: on-demand quadrupedal
assistance for people with limited mobility,'' \emph{J. NeuroEng.
Rehabil.}, vol. 23, Art. no. 180, 2026.

\bibitem{holospot} P. Soler Garcia \emph{et al.}, ``HoloSpot:
Intuitive object manipulation via mixed reality drag-and-drop,''
arXiv:2410.11110, 2024.

\bibitem{miel} A. Oyama, S. Hasegawa, A. Taniguchi, Y. Hagiwara, and
T. Taniguchi, ``Take that for me: Multimodal exophora resolution with
interactive questioning for ambiguous out-of-view instructions,'' in
\emph{Proc. IEEE RO-MAN}, 2025, pp. 563--570.

\bibitem{wang2018} H. Wang, J. Pi, T. Qin, S. Shen, and B. E. Shi,
``SLAM-based localization of 3D gaze using a mobile eye tracker,'' in
\emph{Proc. ACM ETRA}, 2018, Art. no. 65.

\bibitem{attention4d} S. Oishi, K. Koide, M. Yokozuka, and A. Banno,
``4D Attention: Comprehensive framework for spatio-temporal gaze
mapping,'' \emph{IEEE Robot. Autom. Lett.}, vol. 6, no. 4,
pp. 7240--7247, 2021.

\bibitem{semanticfovea} M. Li, N. Songur, P. Orlov, S. Leutenegger,
and A. A. Faisal, ``Towards an embodied semantic fovea: Semantic 3D
scene reconstruction from ego-centric eye-tracker videos,'' presented
at the \emph{ECCV Workshop on Egocentric Perception, Interaction and
Computing (EPIC)}, 2018.

\bibitem{kothari2020} R. Kothari, Z. Yang, C. Kanan, R. Bailey,
J. B. Pelz, and G. J. Diaz, ``Gaze-in-wild: A dataset for studying
eye and head coordination in everyday activities,'' \emph{Sci. Rep.},
vol. 10, Art. no. 2539, 2020.

\end{thebibliography}




\end{document}
